\begin{problem}[Parametric Simultaneous Equations and Matrices]
Consider the simultaneous linear equations involving real constants $m$ and $n$:
\begin{align*}
mx + 2y &= 4 \\
8x + my &= n
\end{align*}

\begin{enumerate}[label=\textbf{(\alph*)}]
    \item Express this system as a matrix equation of the form $AX = B$. Hence, determine the inverse matrix $A^{-1}$ in terms of $m$, stating the specific values of $m$ for which $A^{-1}$ is undefined.
    \item Determine the values of $m$ for which the system does \textit{not} have a unique solution.
    \item Given that $m$ takes the positive value found in part (b), determine the exact value of $n$ that results in an infinite set of solutions. Furthermore, provide a geometric interpretation of the system if $n = 10$ instead.
    \item Suppose $n = 2m$. Assuming a unique solution exists, use the matrix inverse method to solve for $x$ and $y$ in terms of $m$, expressing your final answers in their simplest form.
\end{enumerate}
\end{problem}

\begin{hint}
\begin{itemize}
    \item \textbf{Part (a):} Set up the coefficients of $x$ and $y$ into a $2 \times 2$ matrix. The formula for the inverse includes dividing by the determinant, $ad - bc$. The inverse is undefined when this determinant is equal to zero.
    \item \textbf{Part (b):} A unique solution exists if and only if the coefficient matrix $A$ is invertible. Therefore, a unique solution does \textit{not} exist when the determinant you found in part (a) is zero.
    \item \textbf{Part (c):} Substitute the positive value of $m$ into the original equations. For infinite solutions, the two line equations must be mathematically equivalent (one is a scalar multiple of the other). If the lines have the same gradient but different y-intercepts, what does that mean geometrically?
    \item \textbf{Part (d):} Substitute $n = 2m$ into your column matrix $B$. Multiply $A^{-1}$ by $B$ using standard row-by-column multiplication. Simplify the resulting algebraic expressions for $x$ and $y$ by factoring out common terms.
\end{itemize}
\end{hint}

\begin{solution}
\textbf{(a)} \\
Matrix equation: $\begin{bmatrix} m & 2 \\ 8 & m \end{bmatrix} \begin{bmatrix} x \\ y \end{bmatrix} = \begin{bmatrix} 4 \\ n \end{bmatrix}$. \\
The determinant is $\det(A) = (m)(m) - (2)(8) = m^2 - 16$. \\
$A^{-1} = \frac{1}{m^2 - 16} \begin{bmatrix} m & -2 \\ -8 & m \end{bmatrix}$. \\
The inverse is undefined when $m^2 - 16 = 0$, which occurs when $m = 4$ or $m = -4$.

\medskip
\textbf{(b)} \\
The system does not have a unique solution when the coefficient matrix is singular (i.e., $\det(A) = 0$). \\
From part (a), this occurs when $m = 4$ or $m = -4$.

\medskip
\textbf{(c)} \\
Let $m = 4$. The system becomes: \\
(1) $4x + 2y = 4 \implies 2x + y = 2$ \\
(2) $8x + 4y = n \implies 2x + y = \frac{n}{4}$ \\
For an infinite set of solutions, the equations must represent the same line. Therefore, $2 = \frac{n}{4} \implies n = 8$. \\
If $n = 10$, the equations are $2x + y = 2$ and $2x + y = 2.5$. These equations represent distinct, parallel lines that never intersect, meaning there is no solution.

\medskip
\textbf{(d)} \\
Let $n = 2m$. The matrix $B$ becomes $\begin{bmatrix} 4 \\ 2m \end{bmatrix}$. \\
$X = A^{-1}B = \frac{1}{m^2 - 16} \begin{bmatrix} m & -2 \\ -8 & m \end{bmatrix} \begin{bmatrix} 4 \\ 2m \end{bmatrix}$ \\
$X = \frac{1}{m^2 - 16} \begin{bmatrix} (m)(4) + (-2)(2m) \\ (-8)(4) + (m)(2m) \end{bmatrix}$ \\
$X = \frac{1}{m^2 - 16} \begin{bmatrix} 4m - 4m \\ -32 + 2m^2 \end{bmatrix} = \frac{1}{m^2 - 16} \begin{bmatrix} 0 \\ 2(m^2 - 16) \end{bmatrix}$ \\
$X = \begin{bmatrix} 0 \\ 2 \end{bmatrix}$. \\
The unique solution simplifies elegantly to $x = 0$ and $y = 2$ (for $m \neq \pm 4$).
\end{solution}

\begin{takeaways}
\item \textbf{Parametric Determinants:} Using parameters in the coefficient matrix directly links algebraic singularity (determinant = 0) to geometric relationships (parallel or coincident lines).
\item \textbf{Infinite vs. No Solution:} Once a matrix is singular, comparing the constant terms (the $B$ matrix) dictates whether the lines are coincident (infinite solutions) or strictly parallel (no solution).
\item \textbf{Power of Matrix Algebra:} Part (d) demonstrates how matrix inverse multiplication can elegantly simplify complex parametric systems into clean integer solutions through algebraic factorization.
\end{takeaways}


\begin{problem}[Matrix Systems and Geometric Transformations]
Consider the following system of linear equations:
\begin{align*}
\sqrt{3}x - y &= 4\sqrt{3} \\
x + \sqrt{3}y &= 4
\end{align*}

\begin{enumerate}[label=\textbf{(\alph*)}]
    \item Write this system in the matrix form $AX = B$, and determine the exact value of $\det(A)$.
    \item Determine $A^{-1}$, and hence solve the system of equations algebraically for $X = \begin{bmatrix} x \\ y \end{bmatrix}$.
    \item Express the coefficient matrix $A$ in the form $k \begin{bmatrix} \cos\theta & -\sin\theta \\ \sin\theta & \cos\theta \end{bmatrix}$, where $k > 0$ and $0 \le \theta \le \frac{\pi}{2}$. Hence, describe the two sequential geometric transformations represented by the matrix $A$.
    \item The matrix equation $AX = B$ can be interpreted geometrically as finding a pre-image point $(x, y)$ that maps to the image point $(4\sqrt{3}, 4)$ under transformation $A$. Using your geometric description from part (c), deduce the coordinates of $(x, y)$ using polar coordinates and transformations, verifying your algebraic answer from part (b).
\end{enumerate}
\end{problem}

\begin{hint}
\begin{itemize}
    \item \textbf{Part (a):} Extract the coefficients of $x$ and $y$ carefully. The determinant is calculated as $ad - bc$. Be mindful of squaring surds and negative signs.
    \item \textbf{Part (b):} The inverse of a $2 \times 2$ matrix is $\frac{1}{\det(A)} \begin{bmatrix} d & -b \\ -c & a \end{bmatrix}$. To solve for $X$, pre-multiply the constant matrix $B$ by $A^{-1}$.
    \item \textbf{Part (c):} Notice that the determinant of a pure rotation matrix is always $1$. The determinant of $A$ is not $1$, which implies a dilation is present. The dilation factor $k$ is equal to $\sqrt{\det(A)}$. Factor this $k$ out of matrix $A$ to reveal the standard rotation matrix.
    \item \textbf{Part (d):} Convert the image point $(4\sqrt{3}, 4)$ into polar form $(r, \phi)$. Since applying $A$ dilates and rotates the pre-image \textit{forward}, applying $A^{-1}$ to find the pre-image must do the exact opposite: divide the radius by $k$ and rotate \textit{backwards} by $\theta$.
\end{itemize}
\end{hint}

\begin{solution}
\textbf{(a)} \\
Matrix form: $\begin{bmatrix} \sqrt{3} & -1 \\ 1 & \sqrt{3} \end{bmatrix} \begin{bmatrix} x \\ y \end{bmatrix} = \begin{bmatrix} 4\sqrt{3} \\ 4 \end{bmatrix}$ \\
$\det(A) = (\sqrt{3})(\sqrt{3}) - (-1)(1) = 3 - (-1) = 4$.

\medskip
\textbf{(b)} \\
$A^{-1} = \frac{1}{4}\begin{bmatrix} \sqrt{3} & 1 \\ -1 & \sqrt{3} \end{bmatrix}$. \\
$X = A^{-1}B = \frac{1}{4}\begin{bmatrix} \sqrt{3} & 1 \\ -1 & \sqrt{3} \end{bmatrix}\begin{bmatrix} 4\sqrt{3} \\ 4 \end{bmatrix}$ \\
$X = \frac{1}{4}\begin{bmatrix} \sqrt{3}(4\sqrt{3}) + 1(4) \\ -1(4\sqrt{3}) + \sqrt{3}(4) \end{bmatrix} = \frac{1}{4}\begin{bmatrix} 12 + 4 \\ -4\sqrt{3} + 4\sqrt{3} \end{bmatrix} = \frac{1}{4}\begin{bmatrix} 16 \\ 0 \end{bmatrix} = \begin{bmatrix} 4 \\ 0 \end{bmatrix}$. \\
Therefore, $x = 4$ and $y = 0$.

\medskip
\textbf{(c)} \\
Since $\det(A) = 4$, the dilation factor is $k = \sqrt{4} = 2$. \\
Factoring out $2$: 
$A = 2 \begin{bmatrix} \frac{\sqrt{3}}{2} & -\frac{1}{2} \\ \frac{1}{2} & \frac{\sqrt{3}}{2} \end{bmatrix}$. \\
We need $\cos\theta = \frac{\sqrt{3}}{2}$ and $\sin\theta = \frac{1}{2}$, which gives $\theta = \frac{\pi}{6}$ (or $30^\circ$). \\
Thus, $A = 2 \begin{bmatrix} \cos(\frac{\pi}{6}) & -\sin(\frac{\pi}{6}) \\ \sin(\frac{\pi}{6}) & \cos(\frac{\pi}{6}) \end{bmatrix}$. \\
Geometrically, matrix $A$ represents a \textbf{dilation by a factor of 2} from the origin, followed by an \textbf{anticlockwise rotation of $\frac{\pi}{6}$ radians} about the origin.

\medskip
\textbf{(d)} \\
The image point is $(4\sqrt{3}, 4)$. In polar coordinates: \\
$r = \sqrt{(4\sqrt{3})^2 + 4^2} = \sqrt{48 + 16} = \sqrt{64} = 8$. \\
Angle $\phi = \tan^{-1}\left(\frac{4}{4\sqrt{3}}\right) = \tan^{-1}\left(\frac{1}{\sqrt{3}}\right) = \frac{\pi}{6}$. \\
So the image point is $(8, \frac{\pi}{6})$ in polar. \\
The inverse transformation $A^{-1}$ reverses the effects of $A$: it dilates by $\frac{1}{2}$ and rotates clockwise by $\frac{\pi}{6}$. \\
New radius $r' = 8 \times \frac{1}{2} = 4$. \\
New angle $\phi' = \frac{\pi}{6} - \frac{\pi}{6} = 0$. \\
The pre-image in polar coordinates is $(4, 0)$, which corresponds to the Cartesian coordinates $(4, 0)$. This verifies the algebraic solution $x=4, y=0$.
\end{solution}

\begin{takeaways}
\item \textbf{Algebra Meets Geometry:} Solving a system of linear equations is algebraically identical to finding the pre-image of a point under a 2D matrix transformation.
\item \textbf{Identifying Transformations:} Any matrix of the form $\begin{bmatrix} a & -b \\ b & a \end{bmatrix}$ represents a combined rotation and dilation. The dilation factor is $\sqrt{a^2 + b^2}$ (or $\sqrt{\det(A)}$), and the rotation angle $\theta$ is found using right-triangle trigonometry with $a$ and $b$.
\item \textbf{Inverse Transformations:} If matrix $A$ scales by $k$ and rotates by $\theta$, its inverse $A^{-1}$ scales by $\frac{1}{k}$ and rotates by $-\theta$.
\end{takeaways}


\begin{problem}[Matrix Determinants and System Solutions]
Consider the following system of simultaneous linear equations involving the real parameters $m$ and $p$:
\begin{align*}
mx + 2y &= p \\
12x + (m - 5)y &= 6
\end{align*}

\begin{enumerate}[label=\textbf{(\alph*)}]
    \item Express this system in the matrix form $AX = B$. Determine the values of $m$ for which matrix $A$ is not an invertible matrix.
    \item Let $m$ take the positive value found in part (a). Determine the exact value of $p$ required for the system to have an infinite number of solution pairs. Describe the geometric relationship between the two linear equations in this specific case.
    \item Let $m$ take the negative value found in part (a). Determine the restriction on $p$ such that the system of equations has no solutions. Describe the geometric relationship between the two linear equations in this case.
    \item Suppose $p = 4$, and assume that $m$ is chosen such that a unique solution exists. Use the matrix inverse method to solve for the $x$-coordinate of the solution in terms of $m$, expressing your final answer as a simplified rational fraction.
\end{enumerate}
\end{problem}

\begin{hint}
\begin{itemize}
    \item \textbf{Part (a):} A matrix is not invertible (it is singular) when its determinant equals zero. Extract the coefficients, calculate the determinant using $ad - bc$, set it to zero, and factor the resulting quadratic equation.
    \item \textbf{Part (b):} Substitute your positive $m$ value into the original system. For a linear system to have infinite solutions, the two equations must be scalar multiples of one another (representing coincident lines).
    \item \textbf{Part (c):} Substitute your negative $m$ value into the original system. For a linear system to have no solutions, the lines must have identical gradients but different $y$-intercepts. What value must $p$ \textit{not} equal to prevent the lines from overlapping?
    \item \textbf{Part (d):} The inverse matrix is $A^{-1} = \frac{1}{\det(A)} \begin{bmatrix} d & -b \\ -c & a \end{bmatrix}$. To find $x$, you only need to calculate the top row of the matrix multiplication $A^{-1}B$. Look for common factors in the numerator and denominator to simplify.
\end{itemize}
\end{hint}

\begin{solution}
\textbf{(a)} \\
Matrix equation: $\begin{bmatrix} m & 2 \\ 12 & m-5 \end{bmatrix} \begin{bmatrix} x \\ y \end{bmatrix} = \begin{bmatrix} p \\ 6 \end{bmatrix}$. \\
Matrix $A$ is not invertible when $\det(A) = 0$. \\
$\det(A) = (m)(m-5) - (2)(12) = m^2 - 5m - 24$. \\
Setting to zero: $m^2 - 5m - 24 = 0 \implies (m-8)(m+3) = 0$. \\
Therefore, $A$ is singular when $m = 8$ or $m = -3$.

\medskip
\textbf{(b)} \\
Let $m = 8$. The system becomes: \\
(1) $8x + 2y = p \implies 4x + y = \frac{p}{2}$ \\
(2) $12x + 3y = 6 \implies 4x + y = 2$ \\
For infinite solutions, the equations must be equivalent. \\
Therefore, $\frac{p}{2} = 2 \implies p = 4$. \\
Geometrically, the equations represent \textbf{coincident (identical) straight lines}.

\medskip
\textbf{(c)} \\
Let $m = -3$. The system becomes: \\
(1) $-3x + 2y = p$ \\
(2) $12x - 8y = 6 \implies -3x + 2y = -1.5$ \\
For the system to have no solutions, the equations must be mathematically inconsistent. \\
Therefore, $p$ can be any real number except $-1.5$ (i.e., $p \neq -1.5$). \\
Geometrically, the equations represent \textbf{distinct parallel lines} that never intersect.

\medskip
\textbf{(d)} \\
Let $p = 4$. For a unique solution, $\det(A) = (m-8)(m+3) \neq 0$. \\
$X = A^{-1}B = \frac{1}{(m-8)(m+3)} \begin{bmatrix} m-5 & -2 \\ -12 & m \end{bmatrix} \begin{bmatrix} 4 \\ 6 \end{bmatrix}$. \\
We only need to evaluate the top row for $x$: \\
$x = \frac{4(m-5) + 6(-2)}{(m-8)(m+3)} = \frac{4m - 20 - 12}{(m-8)(m+3)} = \frac{4m - 32}{(m-8)(m+3)}$. \\
Factoring the numerator: \\
$x = \frac{4(m-8)}{(m-8)(m+3)}$. \\
Since $m \neq 8$, we can cancel the common factor, yielding the simplified solution: \\
$x = \frac{4}{m+3}$.
\end{solution}

\begin{takeaways}
\item \textbf{The Zero Determinant:} When $\det(A) = 0$, a unique solution does not exist. This algebraically signals that the lines are either parallel or coincident.
\item \textbf{The Role of Constants:} Once singularity is established by the coefficient matrix $A$, the constant matrix $B$ dictates whether the system collapses into an infinite solution set or an empty set (no solutions).
\item \textbf{Parametric Simplification:} In matrix algebra, finding a unique parametric solution often leads to rational expressions where the roots of the determinant (e.g., $m-8$) appear as cancellable factors in the numerators.
\end{takeaways}


\begin{problem}[Matrix Systems and Parametric Constants]
Consider the following system of linear equations involving a real constant $k$:
\begin{align*}
kx + 2y &= k+1 \\
8x + ky &= 6
\end{align*}

\begin{enumerate}[label=\textbf{(\alph*)}]
    \item Write this system in the matrix form, as $AX = K$. Determine the values of $k$ for which the coefficient matrix $A$ is not an invertible matrix.
    \item For the negative value of $k$ found in part (a), determine the geometric relationship between the two linear equations. How many pairs does the solution set contain?
    \item For the positive value of $k$ found in part (a), demonstrate algebraically whether any solutions can be found for this system of equations. Interpret this result geometrically.
    \item Let $k = 2$. Calculate $A^{-1}$ and hence find the unique solution pair $(x, y)$ that satisfies the system.
\end{enumerate}
\end{problem}

\begin{hint}
\begin{itemize}
    \item \textbf{Part (a):} Matrix $A$ contains the coefficients of $x$ and $y$. It is not invertible (singular) when its determinant, calculated as $ad-bc$, is equal to zero.
    \item \textbf{Part (b):} Substitute the negative $k$ value back into the original equations. If one equation simplifies to become a direct multiple of the other, they represent the exact same line.
    \item \textbf{Part (c):} Substitute the positive $k$ value. Try to solve the system using elimination; you will reach a false mathematical statement, indicating parallel lines that never cross.
    \item \textbf{Part (d):} Use the inverse matrix formula $A^{-1} = \frac{1}{\det(A)} \begin{bmatrix} d & -b \\ -c & a \end{bmatrix}$ and multiply it by the constant column matrix $K$.
\end{itemize}
\end{hint}

\begin{solution}
\textbf{(a)} \\
Matrix form: $\begin{bmatrix} k & 2 \\ 8 & k \end{bmatrix} \begin{bmatrix} x \\ y \end{bmatrix} = \begin{bmatrix} k+1 \\ 6 \end{bmatrix}$. \\
The matrix $A$ is not invertible when $\det(A) = 0$. \\
$\det(A) = (k)(k) - (2)(8) = k^2 - 16$. \\
Setting to zero: $k^2 - 16 = 0 \implies k = 4$ or $k = -4$.

\medskip
\textbf{(b)} \\
Let $k = -4$. The system becomes: \\
(1) $-4x + 2y = -3$ \\
(2) $8x - 4y = 6 \implies -4x + 2y = -3$ \\
The equations are mathematically identical. Geometrically, they represent \textbf{coincident lines}. Because they overlap entirely, the solution set contains an \textbf{infinite number of pairs}.

\medskip
\textbf{(c)} \\
Let $k = 4$. The system becomes: \\
(1) $4x + 2y = 5$ \\
(2) $8x + 4y = 6 \implies 4x + 2y = 3$ \\
Subtracting the simplified second equation from the first gives $0 = 2$, which is a contradiction. Therefore, \textbf{no solutions can be found}. Geometrically, they represent \textbf{distinct parallel lines} that never intersect.

\medskip
\textbf{(d)} \\
For $k = 2$, $A = \begin{bmatrix} 2 & 2 \\ 8 & 2 \end{bmatrix}$ and $K = \begin{bmatrix} 3 \\ 6 \end{bmatrix}$. \\
$\det(A) = (2)(2) - (2)(8) = 4 - 16 = -12$. \\
$A^{-1} = -\frac{1}{12} \begin{bmatrix} 2 & -2 \\ -8 & 2 \end{bmatrix}$. \\
$X = A^{-1}K = -\frac{1}{12} \begin{bmatrix} 2 & -2 \\ -8 & 2 \end{bmatrix} \begin{bmatrix} 3 \\ 6 \end{bmatrix} = -\frac{1}{12} \begin{bmatrix} 6 - 12 \\ -24 + 12 \end{bmatrix} = -\frac{1}{12} \begin{bmatrix} -6 \\ -12 \end{bmatrix} = \begin{bmatrix} \frac{1}{2} \\ 1 \end{bmatrix}$. \\
The unique solution pair is $(\frac{1}{2}, 1)$.
\end{solution}

\begin{takeaways}
\item \textbf{Singularity and Geometry:} A non-invertible (singular) coefficient matrix indicates that the system does not have a unique solution; the lines are either parallel or coincident.
\item \textbf{The Role of Constant Matrices:} The constant terms determine the final state of a singular system. Equivalent ratios across all terms yield infinite solutions (coincident), while differing constant ratios yield zero solutions (parallel).
\item \textbf{Direct Solutions:} The matrix inverse method provides a direct, algorithmic approach to solving uniquely determined systems without requiring manual substitution or elimination.
\end{takeaways}


\begin{problem}[Matrix Algebra and Transformations]
Consider the $3 \times 3$ matrix $A = \begin{bmatrix} 2 & 0 & 0 \\ 0 & 1 & \sqrt{3} \\ 0 & \sqrt{3} & -1 \end{bmatrix}$.

\begin{enumerate}[label=\textbf{(\alph*)}]
    \item Determine $A^2$ and express your result in the form $kI$, where $I$ is the $3 \times 3$ identity matrix and $k$ is a real constant.
    \item Hence, algebraically deduce the inverse matrix $A^{-1}$ without calculating the determinant.
    \item A 3D linear transformation represented by matrix $A$ maps a coordinate point $P(x, y, z)$ to the image point $P'(4, -2, 2\sqrt{3})$. Use your result from part (b) to determine the exact coordinates of the original point $P$.
    \item Using the properties established in part (a), evaluate $A^3$ and hence determine an expression for $A^5$ strictly in terms of $A$.
\end{enumerate}
\end{problem}

\begin{hint}
\begin{itemize}
    \item \textbf{Part (a):} Perform standard $3 \times 3$ matrix multiplication ($A \times A$). Be careful when multiplying the surds in the second and third rows.
    \item \textbf{Part (b):} The definition of an inverse matrix is $A A^{-1} = I$. Take your equation from part (a), $A^2 = kI$, and algebraically manipulate it to isolate $I$ on one side of the equation.
    \item \textbf{Part (c):} The transformation is defined by the equation $AP = P'$. To solve for the pre-image $P$, you must pre-multiply both sides of the equation by $A^{-1}$. 
    \item \textbf{Part (d):} Do not calculate the matrices manually. Use exponent rules for matrices: $A^3 = A^2 \times A$. Substitute the identity relationship you found in part (a) to simplify the powers.
\end{itemize}
\end{hint}

\begin{solution}
\textbf{(a)} \\
$A^2 = \begin{bmatrix} 2 & 0 & 0 \\ 0 & 1 & \sqrt{3} \\ 0 & \sqrt{3} & -1 \end{bmatrix} \begin{bmatrix} 2 & 0 & 0 \\ 0 & 1 & \sqrt{3} \\ 0 & \sqrt{3} & -1 \end{bmatrix}$ \\
$A^2 = \begin{bmatrix} (2)(2)+0+0 & 0+0+0 & 0+0+0 \\ 0+0+0 & (1)(1)+(\sqrt{3})(\sqrt{3})+0 & (1)(\sqrt{3})+(\sqrt{3})(-1)+0 \\ 0+0+0 & (\sqrt{3})(1)+(-1)(\sqrt{3})+0 & (\sqrt{3})(\sqrt{3})+(-1)(-1)+0 \end{bmatrix}$ \\
$A^2 = \begin{bmatrix} 4 & 0 & 0 \\ 0 & 1+3 & \sqrt{3}-\sqrt{3} \\ 0 & \sqrt{3}-\sqrt{3} & 3+1 \end{bmatrix} = \begin{bmatrix} 4 & 0 & 0 \\ 0 & 4 & 0 \\ 0 & 0 & 4 \end{bmatrix}$ \\
$A^2 = 4 \begin{bmatrix} 1 & 0 & 0 \\ 0 & 1 & 0 \\ 0 & 0 & 1 \end{bmatrix} = 4I$. \\
Therefore, $k = 4$.

\medskip
\textbf{(b)} \\
From part (a), $A^2 = 4I$. \\
$A \cdot A = 4I$ \\
Dividing both sides by the scalar $4$: \\
$A \cdot \left(\frac{1}{4}A\right) = I$ \\
Since $A \cdot A^{-1} = I$, it follows that $A^{-1} = \frac{1}{4}A$. \\
$A^{-1} = \begin{bmatrix} 1/2 & 0 & 0 \\ 0 & 1/4 & \sqrt{3}/4 \\ 0 & \sqrt{3}/4 & -1/4 \end{bmatrix}$.

\medskip
\textbf{(c)} \\
Given $AP = P'$, we have $P = A^{-1}P'$. \\
Using $A^{-1} = \frac{1}{4}A$: \\
$P = \frac{1}{4} \begin{bmatrix} 2 & 0 & 0 \\ 0 & 1 & \sqrt{3} \\ 0 & \sqrt{3} & -1 \end{bmatrix} \begin{bmatrix} 4 \\ -2 \\ 2\sqrt{3} \end{bmatrix}$ \\
$P = \frac{1}{4} \begin{bmatrix} (2)(4) + 0 + 0 \\ 0 + (1)(-2) + (\sqrt{3})(2\sqrt{3}) \\ 0 + (\sqrt{3})(-2) + (-1)(2\sqrt{3}) \end{bmatrix}$ \\
$P = \frac{1}{4} \begin{bmatrix} 8 \\ -2 + 6 \\ -2\sqrt{3} - 2\sqrt{3} \end{bmatrix} = \frac{1}{4} \begin{bmatrix} 8 \\ 4 \\ -4\sqrt{3} \end{bmatrix} = \begin{bmatrix} 2 \\ 1 \\ -\sqrt{3} \end{bmatrix}$. \\
The coordinates of $P$ are $(2, 1, -\sqrt{3})$.

\medskip
\textbf{(d)} \\
$A^3 = A^2 \cdot A = (4I)A = 4A$. \\
$A^5 = A^3 \cdot A^2 = (4A)(4I) = 16A$. \\
Alternatively, $A^5 = A(A^2)^2 = A(4I)^2 = 16A$.
\end{solution}

\begin{takeaways}
\item \textbf{The ``Hence'' Command:} In mathematics, the word ``hence'' dictates that you must use the result of the previous section. Finding the inverse via $A^2 = kI$ is often faster and less prone to arithmetic error than manually calculating a $3 \times 3$ determinant and adjugate matrix.
\item \textbf{Matrix Power Cycles:} Matrices that satisfy equations like $A^2 = kI$ act similarly to cyclic numbers in modular arithmetic. High powers of the matrix can be drastically simplified without performing tedious sequential multiplication.
\item \textbf{Geometric Inverses:} Algebraically isolating a coordinate vector by multiplying by an inverse matrix geometrically represents reversing the original 3D transformation to find the starting position.
\end{takeaways}


\begin{problem}[Systems of Linear Equations and 3D Geometry]
Consider the following system of simultaneous linear equations representing three planes in 3D space, where $k$ is a real constant:
\begin{align*}
x + 2y - 3z &= 4 \\
x + y + z &= 6 \\
2x + 3y + kz &= 10
\end{align*}

\begin{enumerate}[label=\textbf{(\alph*)}]
    \item Express this system in the matrix form $AX = B$. Calculate the determinant of the $3 \times 3$ coefficient matrix $A$ in terms of $k$, and hence determine the value of $k$ for which the matrix is singular.
    \item For the value of $k$ found in part (a), demonstrate algebraically that the system is consistent (has solutions) rather than inconsistent (no solutions).
    \item Given that the system is consistent and has infinitely many solutions, let $z = \lambda$. Solve the equations to express $x$ and $y$ in terms of the parameter $\lambda$.
    \item Express the general solution from part (c) as a vector equation of a line in the form $\mathbf{r} = \mathbf{a} + \lambda\mathbf{d}$, and briefly describe the geometric arrangement of the three planes.
\end{enumerate}
\end{problem}

\begin{hint}
\begin{itemize}
    \item \textbf{Part (a):} The determinant of a $3 \times 3$ matrix can be found by expanding across the top row. A matrix is singular when its determinant is equal to zero.
    \item \textbf{Part (b):} Look for a linear combination of the first two equations that perfectly yields the third equation. If the left-hand side and right-hand side both match, the system is consistent.
    \item \textbf{Part (c):} Subtract the second equation from the first to eliminate $x$ and find $y$ in terms of $z$. Then substitute this back into the second equation to find $x$. Finally, replace $z$ with $\lambda$.
    \item \textbf{Part (d):} Separate the constant terms from the terms containing $\lambda$ to form your position vector $\mathbf{a}$ and direction vector $\mathbf{d}$. 
\end{itemize}
\end{hint}

\begin{solution}
\textbf{(a)} \\
Matrix form: $\begin{bmatrix} 1 & 2 & -3 \\ 1 & 1 & 1 \\ 2 & 3 & k \end{bmatrix} \begin{bmatrix} x \\ y \\ z \end{bmatrix} = \begin{bmatrix} 4 \\ 6 \\ 10 \end{bmatrix}$. \\
$\det(A) = 1(1(k) - 1(3)) - 2(1(k) - 1(2)) + (-3)(1(3) - 1(2))$ \\
$\det(A) = (k - 3) - 2(k - 2) - 3(1)$ \\
$\det(A) = k - 3 - 2k + 4 - 3 = -k - 2$. \\
For $A$ to be singular, $\det(A) = 0 \implies -k - 2 = 0 \implies k = -2$.

\medskip
\textbf{(b)} \\
Substitute $k = -2$ into the third equation: $2x + 3y - 2z = 10$. \\
Let Equation 1 be $x + 2y - 3z = 4$ and Equation 2 be $x + y + z = 6$. \\
Adding Equation 1 and Equation 2 gives: \\
$(x + x) + (2y + y) + (-3z + z) = 4 + 6$ \\
$2x + 3y - 2z = 10$. \\
Since this perfectly matches the third equation, the third plane contains the line of intersection of the first two planes. The system is consistent.

\medskip
\textbf{(c)} \\
Let $z = \lambda$. \\
Subtracting Eq 2 from Eq 1: \\
$(x + 2y - 3z) - (x + y + z) = 4 - 6$ \\
$y - 4z = -2 \implies y = 4z - 2 \implies y = 4\lambda - 2$. \\
Substitute $y = 4\lambda - 2$ into Eq 2: \\
$x + (4\lambda - 2) + \lambda = 6$ \\
$x + 5\lambda - 2 = 6 \implies x = 8 - 5\lambda$.

\medskip
\textbf{(d)} \\
From part (c), we have $x = 8 - 5\lambda$, $y = -2 + 4\lambda$, $z = 0 + 1\lambda$. \\
In vector form: \\
$\begin{bmatrix} x \\ y \\ z \end{bmatrix} = \begin{bmatrix} 8 \\ -2 \\ 0 \end{bmatrix} + \lambda \begin{bmatrix} -5 \\ 4 \\ 1 \end{bmatrix}$. \\
Geometrically, the three planes intersect along a single common line (forming a sheaf of planes).
\end{solution}

\begin{takeaways}
\item \textbf{Underdetermined vs. Singular Systems:} The original $2 \times 3$ system implicitly represented infinite solutions because there were more variables than equations. By adding a dependent third equation, we transition to a $3 \times 3$ singular matrix scenario, a core concept in QCAA Specialist Mathematics.
\item \textbf{Singularity and Consistency:} A determinant of zero only tells you a unique solution does not exist. You must verify algebraically whether the planes are parallel/form a prism (inconsistent) or intersect along a line (consistent).
\item \textbf{Parametric Lines:} The process of setting a free variable to $\lambda$ is exactly how the Cartesian equations of intersecting planes are converted into the vector equation of a 3D line.
\end{takeaways}


\begin{problem}[Augmented Matrices and Parametric Solutions]
Consider the following system of linear equations involving four variables ($x, y, z, w$) and a real constant $k$:
\begin{align*}
x + 2y - z + w &= 5 \\
2x - y + 3z - 2w &= 0 \\
3x + y + 2z - w &= k
\end{align*}

\begin{enumerate}[label=\textbf{(\alph*)}]
    \item Set up the augmented matrix for this system. By applying elementary row operations (specifically, examining $R_1 + R_2$), determine the exact value of $k$ for which this system is consistent (has solutions).
    \item Given that the system is consistent, it will have a doubly-infinite family of solutions. Let $z = \lambda$ and $w = \mu$. Express $x$ and $y$ strictly in terms of the parameters $\lambda$ and $\mu$.
    \item A new constraint is added to the system such that $x - y = 3$. Determine the linear relationship that must exist between the parameters $\lambda$ and $\mu$ to satisfy this constraint.
    \item Using your result from part (c), determine the unique solution for the coordinates $(x, y, z, w)$ when $w = 2$.
\end{enumerate}
\end{problem}

\begin{hint}
\begin{itemize}
    \item \textbf{Part (a):} Write the system as an augmented matrix $[A|B]$. Notice that adding the first two rows together creates a row that has the exact same variable coefficients as the third row. For the system to be consistent, the constant terms must also match.
    \item \textbf{Part (b):} Since the third equation is redundant when consistent, use only the first two equations. Substitute $z = \lambda$ and $w = \mu$, moving them to the right-hand side. Then use simultaneous elimination on $x$ and $y$.
    \item \textbf{Part (c):} Substitute your parametric expressions for $x$ and $y$ from part (b) into the equation $x - y = 3$ and group the like parameter terms together.
    \item \textbf{Part (d):} Substitute $w = \mu = 2$ into your relationship from part (c) to find the value of $\lambda$ (which represents $z$). Then, plug both parameter values back into your expressions for $x$ and $y$.
\end{itemize}
\end{hint}

\begin{solution}
\textbf{(a)} \\
The augmented matrix is: \\
$\left[\begin{array}{cccc|c} 1 & 2 & -1 & 1 & 5 \\ 2 & -1 & 3 & -2 & 0 \\ 3 & 1 & 2 & -1 & k \end{array}\right]$ \\
Performing the row operation $R_1 + R_2 \rightarrow R_{new}$: \\
$(1+2)x + (2-1)y + (-1+3)z + (1-2)w = 5 + 0 \implies 3x + y + 2z - w = 5$. \\
Comparing this to Row 3 ($3x + y + 2z - w = k$), we can see that for the system to be mathematically consistent, $k = 5$.

\medskip
\textbf{(b)} \\
Let $z = \lambda$ and $w = \mu$. The system reduces to the first two equations: \\
(1) $x + 2y = 5 + \lambda - \mu$ \\
(2) $2x - y = -3\lambda + 2\mu$ \\
Multiply (2) by 2: $4x - 2y = -6\lambda + 4\mu$. \\
Add to (1): $5x = 5 - 5\lambda + 3\mu \implies x = 1 - \lambda + \frac{3}{5}\mu$. \\
Multiply (1) by 2: $2x + 4y = 10 + 2\lambda - 2\mu$. \\
Subtract (2) from this: $5y = 10 + 5\lambda - 4\mu \implies y = 2 + \lambda - \frac{4}{5}\mu$.

\medskip
\textbf{(c)} \\
Apply the constraint: $x - y = 3$. \\
$\left(1 - \lambda + \frac{3}{5}\mu\right) - \left(2 + \lambda - \frac{4}{5}\mu\right) = 3$ \\
$-1 - 2\lambda + \frac{7}{5}\mu = 3$ \\
$-2\lambda + \frac{7}{5}\mu = 4$. \\
Multiplying by 5 yields the linear relationship: $-10\lambda + 7\mu = 20$.

\medskip
\textbf{(d)} \\
Given $w = 2$, we have $\mu = 2$. \\
Substitute into the relationship from (c): \\
$-10\lambda + 7(2) = 20 \implies -10\lambda + 14 = 20 \implies -10\lambda = 6 \implies \lambda = -\frac{3}{5}$. \\
So, $z = -0.6$. \\
Find $x$: $x = 1 - \left(-\frac{3}{5}\right) + \frac{3}{5}(2) = 1 + \frac{3}{5} + \frac{6}{5} = \frac{14}{5} = 2.8$. \\
Find $y$: $y = 2 + \left(-\frac{3}{5}\right) - \frac{4}{5}(2) = 2 - \frac{3}{5} - \frac{8}{5} = -\frac{1}{5} = -0.2$. \\
The unique solution is $(x, y, z, w) = \left(\frac{14}{5}, -\frac{1}{5}, -\frac{3}{5}, 2\right)$ or $(2.8, -0.2, -0.6, 2)$.
\end{solution}

\begin{takeaways}
\item \textbf{Underdetermined Systems:} Matrices representing systems with fewer independent equations than variables yield infinite families of solutions requiring parameters ($\lambda, \mu$).
\item \textbf{Row Operations for Consistency:} Elementary row operations on augmented matrices quickly reveal inconsistencies (e.g., $0 = c$ where $c \neq 0$). If a row reduces entirely to zero on the left, the right side must also be zero for a solution to exist.
\item \textbf{Applying Constraints:} Adding constraints to a parametric solution family reduces its degrees of freedom. Each new valid equation removes one parameter, eventually leading to a single unique point if enough constraints are provided.
\end{takeaways}


\begin{problem}[Advanced Matrix Products and Inverses]
Let $C = \begin{bmatrix} 2 & -1 & 1 \\ 1 & 0 & 2 \\ -1 & 1 & -1 \end{bmatrix}$ and $D = \begin{bmatrix} -2 & 0 & -2 \\ -1 & -1 & -3 \\ 1 & -1 & 1 \end{bmatrix}$.

\begin{enumerate}[label=\textbf{(\alph*)}]
    \item Evaluate the matrix product $CD$. Hence, determine the inverse matrices $C^{-1}$ and $D^{-1}$ in terms of $C$ and $D$.
    \item A system of simultaneous linear equations is given by:
    \begin{align*}
    -2x - 2z &= 4k \\
    -x - y - 3z &= 2 \\
    x - y + z &= -2k
    \end{align*}
    Express this system in the form $DX = V$, and hence solve for $x, y,$ and $z$ in terms of the real constant $k$.
    \item Determine the specific value of $k$ such that the solution point $(x, y, z)$ lies on the plane defined by $x + z = 4$.
    \item Using the inverse properties established in part (a), solve the matrix equation $C Y D = I$ for the $3 \times 3$ matrix $Y$, where $I$ is the identity matrix.
\end{enumerate}
\end{problem}

\begin{hint}
\begin{itemize}
    \item \textbf{Part (a):} Multiply rows of $C$ by columns of $D$. The result will be a scalar multiple of the identity matrix, $mI$. Use the definition of an inverse ($M M^{-1} = I$) to isolate $I$.
    \item \textbf{Part (b):} The coefficient matrix of the system matches matrix $D$. To solve $DX = V$, pre-multiply both sides by $D^{-1}$. Substitute your expression for $D^{-1}$ from part (a) to avoid recalculating an inverse from scratch.
    \item \textbf{Part (c):} Substitute the expressions for $x$ and $z$ (which you found in terms of $k$ in part b) into the plane equation and solve for $k$.
    \item \textbf{Part (d):} Do not attempt to define $Y$ with 9 unknown variables. Instead, isolate $Y$ algebraically by pre-multiplying by $C^{-1}$ and post-multiplying by $D^{-1}$. Remember that a matrix commutes with its own inverse (e.g., $CD = mI \implies DC = mI$).
\end{itemize}
\end{hint}

\begin{solution}
\textbf{(a)} \\
$CD = \begin{bmatrix} 2 & -1 & 1 \\ 1 & 0 & 2 \\ -1 & 1 & -1 \end{bmatrix} \begin{bmatrix} -2 & 0 & -2 \\ -1 & -1 & -3 \\ 1 & -1 & 1 \end{bmatrix} = \begin{bmatrix} -4+1+1 & 0+1-1 & -4+3+1 \\ -2+0+2 & 0+0-2 & -2+0+2 \\ 2-1-1 & 0-1+1 & 2-3-1 \end{bmatrix}$ \\
$CD = \begin{bmatrix} -2 & 0 & 0 \\ 0 & -2 & 0 \\ 0 & 0 & -2 \end{bmatrix} = -2I$. \\
Since $CD = -2I$, we have $C\left(-\frac{1}{2}D\right) = I$, so $C^{-1} = -\frac{1}{2}D$. \\
Similarly, $\left(-\frac{1}{2}C\right)D = I$, so $D^{-1} = -\frac{1}{2}C$.

\medskip
\textbf{(b)} \\
The system is $D \begin{bmatrix} x \\ y \\ z \end{bmatrix} = \begin{bmatrix} 4k \\ 2 \\ -2k \end{bmatrix}$. \\
$X = D^{-1}V = -\frac{1}{2}C V = -\frac{1}{2} \begin{bmatrix} 2 & -1 & 1 \\ 1 & 0 & 2 \\ -1 & 1 & -1 \end{bmatrix} \begin{bmatrix} 4k \\ 2 \\ -2k \end{bmatrix}$ \\
$X = -\frac{1}{2} \begin{bmatrix} 2(4k) - 1(2) + 1(-2k) \\ 1(4k) + 0(2) + 2(-2k) \\ -1(4k) + 1(2) - 1(-2k) \end{bmatrix} = -\frac{1}{2} \begin{bmatrix} 6k - 2 \\ 0 \\ -2k + 2 \end{bmatrix} = \begin{bmatrix} 1 - 3k \\ 0 \\ k - 1 \end{bmatrix}$. \\
Thus, $x = 1 - 3k$, $y = 0$, and $z = k - 1$.

\medskip
\textbf{(c)} \\
We require $x + z = 4$. \\
Substituting the results from (b): \\
$(1 - 3k) + (k - 1) = 4$ \\
$-2k = 4 \implies k = -2$.

\medskip
\textbf{(d)} \\
Given $CYD = I$. \\
Pre-multiply by $C^{-1}$ and post-multiply by $D^{-1}$: \\
$Y = C^{-1} I D^{-1} = C^{-1} D^{-1}$. \\
Substitute the inverses from part (a): \\
$Y = \left(-\frac{1}{2}D\right)\left(-\frac{1}{2}C\right) = \frac{1}{4}DC$. \\
Since $CD = -2I$, it follows that $DC = -2I$ (a matrix and its inverse commute). \\
$Y = \frac{1}{4}(-2I) = -\frac{1}{2}I = \begin{bmatrix} -0.5 & 0 & 0 \\ 0 & -0.5 & 0 \\ 0 & 0 & -0.5 \end{bmatrix}$.
\end{solution}

\begin{takeaways}
\item \textbf{The ``Hence'' Method:} Calculating $3 \times 3$ inverses from scratch using determinants and adjugates is tedious. Finding a product $AB = mI$ is a standard shortcut to define $A^{-1} = \frac{1}{m}B$.
\item \textbf{Parametric Systems:} Matrices handle parametric constants (like $k$) effortlessly because matrix multiplication distributes linearly over expressions.
\item \textbf{Matrix Algebra vs. Arithmetic:} Part (d) demonstrates the power of abstract matrix algebra. Treating matrices as single algebraic entities ($C^{-1}$, $D^{-1}$) avoids massive blocks of arithmetic calculation.
\end{takeaways}


\begin{problem}[Matrix Systems and Planes in 3D Space]
Consider the points $A(1, 2, -1)$, $B(2, 3, 1)$, and $C_k(3, -1, k)$, where $k$ is a real constant. A plane containing these three points can be modeled by the equation $ax + by + cz = 1$, provided the plane does not pass through the origin.

\begin{enumerate}[label=\textbf{(\alph*)}]
    \item Substitute the coordinates of the three points into the plane equation to formulate a system of three linear equations in terms of $a$, $b$, and $c$. Express this system as a matrix equation $M \mathbf{v} = \mathbf{u}$, and find $\det(M)$ in terms of $k$.
    \item Determine the value of $k$ for which the matrix $M$ is singular. Explain why the matrix equation $M \mathbf{v} = \mathbf{u}$ has no solution for this value of $k$, and what this implies about the geometric position of the plane containing $A$, $B$, and $C_k$ relative to the origin.
    \item Let $k = 2$, corresponding to the specific point $C_2(3, -1, 2)$. Solve the system using Gaussian elimination (row reduction) to find the values of $a, b,$ and $c$, and hence state the unique equation of the plane in the form $ax + by + cz + d = 0$ with integer coefficients.
    \item For the singular case where $k = 18$, the plane must pass through the origin. Set up a new homogeneous matrix equation for the plane $a'x + b'y + c'z = 0$, and find the specific Cartesian equation of this plane by letting $a' = 5$.
\end{enumerate}
\end{problem}

\begin{hint}
\begin{itemize}
    \item \textbf{Part (a):} Replace $x, y,$ and $z$ with the coordinates of each point to create three distinct equations. Extract the coefficients of $a, b,$ and $c$ to form the $3 \times 3$ matrix $M$. Calculate the determinant by expanding across the top row.
    \item \textbf{Part (b):} A matrix is singular when its determinant equals zero. If setting $ax+by+cz=1$ yields a singular matrix and an inconsistent system, it means no such plane exists that misses the origin. 
    \item \textbf{Part (c):} Form the augmented matrix $[M | \mathbf{u}]$ and use elementary row operations to reach upper triangular form. Back-substitute to find $c, b,$ and $a$. Multiply the entire plane equation by the lowest common denominator to remove fractions.
    \item \textbf{Part (d):} A plane through the origin equates to $0$ rather than $1$. Substitute the points $A$, $B$, and $C_{18}$ to form a homogeneous system $M\mathbf{v} = \mathbf{0}$. Use the first two equations to express $b'$ and $c'$ strictly in terms of $a'$, then substitute $a'=5$.
\end{itemize}
\end{hint}

\begin{solution}
\textbf{(a)} \\
Substituting the points into $ax + by + cz = 1$: \\
$A(1, 2, -1) \implies a + 2b - c = 1$ \\
$B(2, 3, 1) \implies 2a + 3b + c = 1$ \\
$C_k(3, -1, k) \implies 3a - b + kc = 1$ \\
Matrix equation: $\begin{bmatrix} 1 & 2 & -1 \\ 2 & 3 & 1 \\ 3 & -1 & k \end{bmatrix} \begin{bmatrix} a \\ b \\ c \end{bmatrix} = \begin{bmatrix} 1 \\ 1 \\ 1 \end{bmatrix}$. \\
$\det(M) = 1(3k - (-1)) - 2(2k - 3) + (-1)(-2 - 9) = 3k + 1 - 4k + 6 + 11 = 18 - k$.

\medskip
\textbf{(b)} \\
Matrix $M$ is singular when $\det(M) = 0 \implies 18 - k = 0 \implies k = 18$. \\
When $k=18$, the system is inconsistent (parallel planes in the row space), meaning no unique solution exists for $ax + by + cz = 1$. Geometrically, this implies the plane containing $A, B$, and $C_{18}$ must pass exactly through the origin, meaning its true equation is of the form $ax + by + cz = 0$, not $1$.

\medskip
\textbf{(c)} \\
Let $k = 2$. The augmented matrix is: \\
$\left[\begin{array}{ccc|c} 1 & 2 & -1 & 1 \\ 2 & 3 & 1 & 1 \\ 3 & -1 & 2 & 1 \end{array}\right]$ \\
$R_2 \to R_2 - 2R_1 \implies \left[\begin{array}{ccc|c} 1 & 2 & -1 & 1 \\ 0 & -1 & 3 & -1 \\ 3 & -1 & 2 & 1 \end{array}\right]$ \\
$R_3 \to R_3 - 3R_1 \implies \left[\begin{array}{ccc|c} 1 & 2 & -1 & 1 \\ 0 & -1 & 3 & -1 \\ 0 & -7 & 5 & -2 \end{array}\right]$ \\
$R_3 \to R_3 - 7R_2 \implies \left[\begin{array}{ccc|c} 1 & 2 & -1 & 1 \\ 0 & -1 & 3 & -1 \\ 0 & 0 & -16 & 5 \end{array}\right]$ \\
From $R_3$: $-16c = 5 \implies c = -\frac{5}{16}$. \\
From $R_2$: $-b + 3(-\frac{5}{16}) = -1 \implies b = \frac{1}{16}$. \\
From $R_1$: $a + 2(\frac{1}{16}) - (-\frac{5}{16}) = 1 \implies a = \frac{9}{16}$. \\
The equation is $\frac{9}{16}x + \frac{1}{16}y - \frac{5}{16}z = 1$, which simplifies to $9x + y - 5z - 16 = 0$.

\medskip
\textbf{(d)} \\
For $k = 18$, the homogeneous system is: \\
(1) $a' + 2b' - c' = 0$ \\
(2) $2a' + 3b' + c' = 0$ \\
(3) $3a' - b' + 18c' = 0$ \\
Add (1) and (2) to eliminate $c'$: $3a' + 5b' = 0 \implies b' = -\frac{3}{5}a'$. \\
Let $a' = 5 \implies b' = -3$. \\
Substitute into (1): $5 + 2(-3) - c' = 0 \implies -1 - c' = 0 \implies c' = -1$. \\
The equation of the plane is $5x - 3y - z = 0$.
\end{solution}

\begin{takeaways}
\item \textbf{Non-Homogeneous vs Homogeneous Form:} Setting a plane's constant $d$ to a non-zero value allows for a quick $3 \times 3$ matrix setup, but it operates on the assumption that the plane does not cross the origin.
\item \textbf{Geometric Meaning of Singularity:} If the resulting matrix determinant is zero, it algebraically signals that the initial assumption (the plane avoids the origin) is false for that specific set of points.
\item \textbf{Gaussian Elimination:} Row reduction is a powerful tool for solving $3 \times 3$ systems without needing to calculate cumbersome matrix inverses or adjugates manually.
\end{takeaways}


\begin{problem}[Matrix Polynomials and System Solutions]
Let $A = \begin{bmatrix} 3 & -2 & -2 \\ 3 & -2 & -3 \\ 1 & -1 & 0 \end{bmatrix}$.

\begin{enumerate}[label=\textbf{(\alph*)}]
    \item Evaluate $A^2$ to show that $A$ is an involutory matrix (i.e., a matrix where $A^2 = I$).
    \item Let $E = \frac{1}{2}(I + A)$. Using matrix algebra and your result from part (a), prove that $E$ is an idempotent matrix (i.e., $E^2 = E$).
    \item Consider the matrix $M = I + kA$, where $k$ is a real constant. By expanding the algebraic product $M(I - kA)$, determine the exact values of $k$ for which the matrix $M$ does not have an inverse.
    \item Let $k = 2$. Use your algebraic result from part (c) to determine $M^{-1}$, and hence solve the following system of linear equations:
    \begin{align*}
    7x - 4y - 4z &= 3 \\
    6x - 3y - 6z &= -3 \\
    2x - 2y + z &= 6
    \end{align*}
\end{enumerate}
\end{problem}

\begin{hint}
\begin{itemize}
    \item \textbf{Part (a):} Multiply $A$ by itself using standard row-by-column multiplication. Verify that the main diagonal consists of $1$s and all other entries are $0$.
    \item \textbf{Part (b):} Square the algebraic expression for $E$. Remember that scalar multiples can be factored out, and matrix multiplication distributes over addition. Substitute $A^2 = I$ when simplifying.
    \item \textbf{Part (c):} Expand $(I + kA)(I - kA)$ as you would a difference of perfect squares, keeping in mind that $I$ commutes with any matrix. For $M$ to be singular, the scalar coefficient of $I$ in your resulting expression must equal zero.
    \item \textbf{Part (d):} Notice that the coefficient matrix of the system is precisely $M$ when $k=2$. Calculate $M^{-1}$ by setting $k=2$ in the identity from part (c), then pre-multiply the constant column vector by $M^{-1}$.
\end{itemize}
\end{hint}

\begin{solution}
\textbf{(a)} \\
$A^2 = \begin{bmatrix} 3 & -2 & -2 \\ 3 & -2 & -3 \\ 1 & -1 & 0 \end{bmatrix} \begin{bmatrix} 3 & -2 & -2 \\ 3 & -2 & -3 \\ 1 & -1 & 0 \end{bmatrix}$ \\
$A^2 = \begin{bmatrix} 9-6-2 & -6+4+2 & -6+6+0 \\ 9-6-3 & -6+4+3 & -6+6+0 \\ 3-3+0 & -2+2+0 & -2+3+0 \end{bmatrix} = \begin{bmatrix} 1 & 0 & 0 \\ 0 & 1 & 0 \\ 0 & 0 & 1 \end{bmatrix} = I$.

\medskip
\textbf{(b)} \\
$E^2 = \left(\frac{1}{2}(I + A)\right)^2 = \frac{1}{4}(I + A)(I + A)$ \\
$E^2 = \frac{1}{4}(I^2 + IA + AI + A^2)$. \\
Since $I^2 = I$, $IA = AI = A$, and $A^2 = I$: \\
$E^2 = \frac{1}{4}(I + A + A + I) = \frac{1}{4}(2I + 2A) = \frac{1}{2}(I + A) = E$.

\medskip
\textbf{(c)} \\
$M(I - kA) = (I + kA)(I - kA) = I^2 - kIA + kAI - k^2 A^2$. \\
Because $IA = AI = A$, the middle terms cancel: \\
$M(I - kA) = I - k^2 A^2$. \\
Substitute $A^2 = I$: \\
$M(I - kA) = I - k^2 I = (1 - k^2)I$. \\
Rearranging this gives $M \left[\frac{1}{1-k^2}(I-kA)\right] = I$. \\
Matrix $M$ is singular (has no inverse) when the scalar divisor $1 - k^2 = 0$. \\
Therefore, $M$ cannot have an inverse when $k = 1$ or $k = -1$.

\medskip
\textbf{(d)} \\
Let $k = 2$. The system's coefficient matrix is $M = I + 2A = \begin{bmatrix} 7 & -4 & -4 \\ 6 & -3 & -6 \\ 2 & -2 & 1 \end{bmatrix}$. \\
From (c), $M^{-1} = \frac{1}{1-2^2}(I - 2A) = -\frac{1}{3}(I - 2A)$. \\
$I - 2A = \begin{bmatrix} 1 & 0 & 0 \\ 0 & 1 & 0 \\ 0 & 0 & 1 \end{bmatrix} - \begin{bmatrix} 6 & -4 & -4 \\ 6 & -4 & -6 \\ 2 & -2 & 0 \end{bmatrix} = \begin{bmatrix} -5 & 4 & 4 \\ -6 & 5 & 6 \\ -2 & 2 & 1 \end{bmatrix}$. \\
$M^{-1} = -\frac{1}{3} \begin{bmatrix} -5 & 4 & 4 \\ -6 & 5 & 6 \\ -2 & 2 & 1 \end{bmatrix}$. \\
$X = M^{-1}B = -\frac{1}{3} \begin{bmatrix} -5 & 4 & 4 \\ -6 & 5 & 6 \\ -2 & 2 & 1 \end{bmatrix} \begin{bmatrix} 3 \\ -3 \\ 6 \end{bmatrix}$ \\
$X = -\frac{1}{3} \begin{bmatrix} -15 - 12 + 24 \\ -18 - 15 + 36 \\ -6 - 6 + 6 \end{bmatrix} = -\frac{1}{3} \begin{bmatrix} -3 \\ 3 \\ -6 \end{bmatrix} = \begin{bmatrix} 1 \\ -1 \\ 2 \end{bmatrix}$. \\
Therefore, $x = 1$, $y = -1$, and $z = 2$.
\end{solution}

\begin{takeaways}
\item \textbf{Abstract Matrix Algebra:} Abstract algebra simplifies complex computational problems. Establishing $A^2=I$ creates a generalized shortcut for inverting any matrix of the parametric form $I+kA$.
\item \textbf{Polynomial Singularity Proofs:} The QCAA curriculum emphasizes proving invertibility using polynomial identities (like $(I+kA)(I-kA) = (1-k^2)I$) rather than relying exclusively on calculating massive $3 \times 3$ determinants.
\item \textbf{Idempotent and Involutory Links:} Idempotent matrices ($E^2=E$, which geometrically represent projections) and involutory matrices ($A^2=I$, which geometrically represent reflections) are fundamentally linked via simple algebraic transformations like $E = \frac{1}{2}(I + A)$.
\end{takeaways}
